%%%%%%%%%%%%%%%%%%%%%%%%%%%%%%%%%%%%%%%%%%%%%%%%%%%%%%%%%%%%%%%%%%%%%%%%%%%%%%%
% APPENDIX DOMAIN-CONSULTING: BEHAVIORAL CONSULTING METHODOLOGY
%%%%%%%%%%%%%%%%%%%%%%%%%%%%%%%%%%%%%%%%%%%%%%%%%%%%%%%%%%%%%%%%%%%%%%%%%%%%%%%
%
% CATEGORY: DOMAIN (Applications)
% BEATRIX Module: BCM2_DOMAIN_CONSULTING_9200
% Version: 1.0 (January 2026)
% Category: DOMAIN
% Language: English
%
% PURPOSE:
% Demonstrates how EBF transforms consulting practice through the Framework-Model
% distinction, quantified comparisons via 20 KPIs, and technology-enabled
% behavioral analysis with Beatrix.
%
% KEY QUESTION: How does EBF-based consulting differ from classical consulting,
% and what are the quantifiable advantages?
%
% LANGUAGE: English
%
%%%%%%%%%%%%%%%%%%%%%%%%%%%%%%%%%%%%%%%%%%%%%%%%%%%%%%%%%%%%%%%%%%%%%%%%%%%%%%%

\chapter{DOMAIN-CONSULTING: Behavioral Consulting Methodology}
\label{app:domain-consulting}

%==============================================================================
% HEADER BLOCK (Template Compliance)
%==============================================================================

\begin{tcolorbox}[colback=blue!5!white, colframe=blue!75!black,
    title=\textbf{Appendix UNMAPPED_BHC: Application Domain}]

\begin{tabular}{@{}ll@{}}
\textbf{Appendix:} & DOMAIN-CONSULTING (Behavioral Consulting Methodology) \\
\textbf{Category:} & DOMAIN (Applications) \\
\textbf{Scope:} & Consulting practice transformation through EBF \\
\textbf{Version:} & 1.0 \\
\textbf{Last Updated:} & January 2026 \\
\textbf{Dependencies:} & Appendix UNMAPPED_MSC (Framework vs Model theory) \\
 & Appendix CTW (Context dimensions, SwissHealth case) \\
 & Appendix UNMAPPED_GLS (Glossary) \\
\textbf{Outputs:} & 10 EBF advantages, 20 KPIs, Implementation roadmap \\
\end{tabular}

\end{tcolorbox}

%==============================================================================
% CROSS-REFERENCE MAP
%==============================================================================

\begin{tcolorbox}[colback=gray!5!white, colframe=gray!75!black,
    title=\textbf{Cross-Reference Map}]

\textbf{This Appendix $\rightarrow$ Other Documents:}
\begin{itemize}[nosep]
    \item Appendix UNMAPPED_MSC: Framework vs Model epistemology (Section~\ref{sec:meta-framework-model})
    \item Appendix CTW: SwissHealth worked example (Section~\ref{sec:V-example-org})
    \item Appendix UNMAPPED_GLS: Central notation and definitions
    \item Appendix LLM1: Beatrix technical architecture
    \item Bibliography: \texttt{bcm\_master.bib}
\end{itemize}

\textbf{Other Documents $\rightarrow$ This Appendix:}
\begin{itemize}[nosep]
    \item Chapter 14: References consulting methodology
    \item Chapter 15: References KPI framework
\end{itemize}

\end{tcolorbox}

%------------------------------------------------------------------------------
% CHAPTER LINKAGE BOX
%------------------------------------------------------------------------------
\begin{tcolorbox}[colback=purple!5!white, colframe=purple!75!black,
    title=Chapter Linkage]

\textbf{Primary Connection:}
\begin{itemize}[nosep]
\item Appendix UNMAPPED_MSC: Theoretical foundation for Framework vs Model
\item Appendix CTW: Practical application (SwissHealth case)
\end{itemize}

\textbf{This Appendix Answers:}
\begin{itemize}[nosep]
\item Why do consulting firms use standardized frameworks?
\item How does EBF differ from classical consulting frameworks?
\item What are the quantifiable advantages (KPIs)?
\item How does Beatrix operationalize the EBF framework?
\end{itemize}

\end{tcolorbox}

%------------------------------------------------------------------------------
% ABSTRACT
%------------------------------------------------------------------------------
\begin{abstract}
This appendix establishes behavioral consulting as a distinct application domain for EBF. We demonstrate how the Framework-Model distinction transforms consulting practice: classical frameworks (McKinsey 7S, Porter's Five Forces, BCG Matrix) are \textit{descriptive templates} that organize analysis but do not generate predictions. EBF is a \textit{generative framework} that produces testable models with quantified predictions. We present 10 structural advantages of EBF over classical approaches and validate these through 20 KPIs measuring both consulting firm performance and client value creation. The Beatrix system operationalizes this framework, enabling technology-assisted behavioral consulting with 3-5$\times$ higher success rates and ROI improvements exceeding 1,000\%.
\end{abstract}

%------------------------------------------------------------------------------
% QUICK REFERENCE
%------------------------------------------------------------------------------
\begin{tcolorbox}[colback=blue!5!white, colframe=blue!75!black,
    title=Quick Reference: EBF Consulting Advantage]

\textbf{The Core Distinction:}
\begin{quote}
Classical frameworks help you \textbf{describe} what is.\\
EBF helps you \textbf{predict} what will be.
\end{quote}

\textbf{Key Metrics:}
\begin{itemize}[nosep]
\item Implementation success: 30\% (classical) $\rightarrow$ 75\% (EBF)
\item Time to diagnosis: 6 weeks $\rightarrow$ 5 days (-85\%)
\item Client ROI: 150\% $\rightarrow$ 450\% (3$\times$)
\item Cost per outcome: -84\%
\end{itemize}

\textbf{Key Insight:} The difference is between HANDICRAFT and INDUSTRIALIZATION.

\end{tcolorbox}

%==============================================================================
% FUNDAMENTAL QUESTION (Template Compliance)
%==============================================================================

\begin{tcolorbox}[colback=red!5!white, colframe=red!75!black,
    title=\textbf{Fundamental Question}]

\textbf{How does EBF-based consulting differ from classical consulting approaches,
and what are the quantifiable advantages?}

\vspace{0.5em}
\textbf{Answer:}
\begin{enumerate}[nosep]
    \item Classical frameworks are \textbf{descriptive templates}---they organize thinking but don't predict outcomes
    \item EBF is a \textbf{generative framework}---it produces testable models with quantified predictions
    \item The advantage is measurable: 10 structural benefits, validated through 20 KPIs
    \item Beatrix operationalizes the complexity that makes EBF superior
\end{enumerate}

\end{tcolorbox}


%%%%%%%%%%%%%%%%%%%%%%%%%%%%%%%%%%%%%%%%%%%%%%%%%%%%%%%%%%%%%%%%%%%%%%%%%%%%%%%
% SECTION 1: THE CONSULTING FRAMEWORK PROBLEM
%%%%%%%%%%%%%%%%%%%%%%%%%%%%%%%%%%%%%%%%%%%%%%%%%%%%%%%%%%%%%%%%%%%%%%%%%%%%%%%

% -----------------------------------------------------------------------------
% APPENDIX SCOPE BOX
% -----------------------------------------------------------------------------

\begin{tcolorbox}[
    colback=orange!5!white,
    colframe=orange!75!black,
    title={\textbf{Appendix Scope: Ziel / In-Scope / Out-of-Scope / Constraints}},
    fonttitle=\bfseries
]

\textbf{Ziel (Objective):}
\begin{itemize}[nosep]
    \item How does EBF-based consulting differ from classical consulting approaches,
and what are the quantifiable advantages?
\end{itemize}

\textbf{In-Scope:}
\begin{itemize}[nosep]
    \item Introduction: The Consulting Framework Problem
    \item Framework vs.\ Model: The Core Distinction
    \item The 10 Advantages of EBF over Classical Frameworks
    \item The 20 KPIs: Quantified Comparison
    \item Beatrix: Technology-Enabled Behavioral Consulting
\end{itemize}

\textbf{Out-of-Scope (delegiert an andere Appendices):}
\begin{itemize}[nosep]
    \item Framework vs Model theory $\rightarrow$ Appendix UNMAPPED_MSC
    \item Context dimensions, SwissHealth case $\rightarrow$ Appendix CTW
    \item Glossary $\rightarrow$ Appendix UNMAPPED_GLS
\end{itemize}

\textbf{Constraints (Anwendungsgrenzen):}
\begin{itemize}[nosep]
    \item [TODO: Define when this appendix does NOT apply]
    \item [TODO: Define required prerequisites]
    \item [TODO: Define known limitations]
\end{itemize}

\textbf{Lieferobjekte (Deliverables):}
\begin{itemize}[nosep]
    \item 5 Table(s)
    \item 3 Equation(s)
    \item Worked Example(s)
\end{itemize}

\end{tcolorbox}


\section{Introduction: The Consulting Framework Problem}
\label{sec:consulting-intro}

\subsection{Why Consultants Use Frameworks}

Every major consulting firm relies on standardized frameworks:

\begin{center}
\renewcommand{\arraystretch}{1.3}
\begin{tabular}{@{}llp{6cm}@{}}
\toprule
\textbf{Framework} & \textbf{Developer} & \textbf{Purpose} \\
\midrule
7S Model & McKinsey & Organizational analysis \\
Five Forces & Porter/Monitor & Industry analysis \\
BCG Matrix & BCG & Portfolio strategy \\
Value Chain & Porter & Competitive advantage \\
SWOT & Various & Situational assessment \\
\bottomrule
\end{tabular}
\end{center}

\textbf{Why frameworks?} Four business reasons:

\begin{enumerate}
\item \textbf{Completeness:} Ensures no dimension is forgotten
\item \textbf{Comparability:} All projects analyze same dimensions $\rightarrow$ cross-learning
\item \textbf{Scalability:} Juniors fill templates, partners interpret
\item \textbf{Quality control:} ``All fields filled'' = analysis complete
\end{enumerate}

\subsection{The Template (``Schablone'') Logic}

\begin{tcolorbox}[colback=gray!5!white, colframe=gray!75!black,
    title=\textbf{Frameworks as Templates}]

A framework is a \textbf{template with empty fields}:

\begin{center}
\begin{tabular}{cc}
\textbf{Template (Framework)} & \textbf{Filled In (Model)} \\
\midrule
\begin{tabular}{|c|c|c|}
\hline
\rule{1cm}{0pt} & \rule{1cm}{0pt} & \rule{1cm}{0pt} \\
\hline
\rule{1cm}{0pt} & \rule{1cm}{0pt} & \rule{1cm}{0pt} \\
\hline
\end{tabular}
&
\begin{tabular}{|c|c|c|}
\hline
A & B & C \\
\hline
D & E & F \\
\hline
\end{tabular}
\\[0.5em]
``These fields must be filled'' & ``For client X, these values apply''
\end{tabular}
\end{center}

\textbf{The template guarantees:}
\begin{itemize}[nosep]
    \item Every relevant dimension is considered
    \item Analysis is systematic, not random
    \item Results are comparable across projects
\end{itemize}

\end{tcolorbox}

\subsection{What Classical Frameworks Miss}

Classical frameworks share a critical limitation:

\begin{tcolorbox}[colback=red!5!white, colframe=red!75!black,
    title=\textbf{The Missing Element}]

\textbf{Classical frameworks are DESCRIPTIVE, not PREDICTIVE.}

\begin{center}
\renewcommand{\arraystretch}{1.3}
\begin{tabular}{@{}p{6cm}p{6cm}@{}}
\textbf{What Five Forces Says:} & \textbf{What It Doesn't Say:} \\
\midrule
``Buyer power is high'' & ``If you do X, buyer power decreases by Y\%'' \\
``Threat of substitutes is medium'' & ``Intervention Z will reduce threat with probability P'' \\
\end{tabular}
\end{center}

\textbf{The consultant must bridge the gap} from description to recommendation
using intuition, experience, and guesswork.

\end{tcolorbox}


%%%%%%%%%%%%%%%%%%%%%%%%%%%%%%%%%%%%%%%%%%%%%%%%%%%%%%%%%%%%%%%%%%%%%%%%%%%%%%%
% SECTION 2: FRAMEWORK VS MODEL
%%%%%%%%%%%%%%%%%%%%%%%%%%%%%%%%%%%%%%%%%%%%%%%%%%%%%%%%%%%%%%%%%%%%%%%%%%%%%%%

\section{Framework vs.\ Model: The Core Distinction}
\label{sec:consulting-fw-model}

\subsection{Definitions for Practitioners}

\begin{tcolorbox}[colback=blue!5!white, colframe=blue!75!black,
    title=\textbf{The Cookbook Analogy}]

\textbf{Framework = Cookbook structure}
\begin{quote}
``Every recipe has: ingredients, preparation, cooking time, serving suggestion''
\end{quote}
This tells you \textit{how to think about cooking}. It doesn't tell you \textit{what to cook}.

\textbf{Model = Specific recipe}
\begin{quote}
``Spaghetti Carbonara: 400g pasta, 150g guanciale, 4 egg yolks, 100g pecorino $\rightarrow$ 12 minutes $\rightarrow$ serves 4''
\end{quote}
This makes a \textit{specific, testable prediction}: these inputs $\rightarrow$ this output.

\end{tcolorbox}

\begin{tcolorbox}[colback=green!5!white, colframe=green!75!black,
    title=\textbf{The Map Analogy}]

\textbf{Framework = Map style} (topographic, political, road map)

\textbf{Model = Specific route} (``Take A1 to Bern, exit 24'')

\vspace{0.5em}
\begin{tabular}{@{}p{6cm}p{6cm}@{}}
\textbf{Map Style (Framework)} & \textbf{Route (Model)} \\
\midrule
Shows \textit{what to look for} & Tells \textit{what to do} \\
Multiple routes possible & One specific prediction \\
Can't be ``wrong''---only useful or not & Can be wrong (road closed) \\
Integrates new roads easily & Fails if road doesn't exist \\
\end{tabular}

\vspace{0.5em}
\textbf{Key insight:} When your route fails (traffic jam), you don't throw away the map.
You use the map to find another route. But economics threw away the map when routes failed.

\end{tcolorbox}

\subsection{The Formal Distinction}

\begin{center}
\renewcommand{\arraystretch}{1.4}
\begin{tabular}{@{}p{3cm}p{5.5cm}p{5.5cm}@{}}
\toprule
& \textbf{Framework} & \textbf{Model} \\
\midrule
\textbf{Purpose} & Organizes thinking & Makes predictions \\
\textbf{Claim} & ``Here is how to think about X'' & ``X equals Y'' \\
\textbf{Falsifiable?} & No (meta-level) & Yes (testable) \\
\textbf{On failure} & Generates next hypothesis & Must be discarded \\
\textbf{Scope} & Broad, organizing & Narrow, specific \\
\textbf{Example} & EBF's 10C $\times$ 8$\Psi$ & ``Weekly rewards increase participation by 12\%'' \\
\bottomrule
\end{tabular}
\end{center}

\subsection{Why This Matters for Consulting}

\begin{tcolorbox}[colback=yellow!5!white, colframe=orange!75!black,
    title=\textbf{Without vs.\ With the Framework-Model Distinction}]

\textbf{Without (Model-Only Thinking):}
\begin{enumerate}[nosep]
    \item Client asks: ``Does nudging work?''
    \item Consultant tests specific nudge
    \item Nudge fails
    \item Conclusion: ``Nudging doesn't work here''
    \item Client hires different consultant
\end{enumerate}

\textbf{With (Framework Thinking):}
\begin{enumerate}[nosep]
    \item Client asks: ``Does nudging work?''
    \item Consultant analyzes via framework (10C, 8$\Psi$)
    \item Identifies: $\Psi_T$ (temporal) is the binding constraint
    \item Designs context intervention targeting $\Psi_T$
    \item If intervention fails, framework guides next attempt
    \item Systematic improvement, not trial-and-error
\end{enumerate}

\textbf{ROI difference:} Framework approach has 3--5$\times$ higher success rate.

\end{tcolorbox}

\textit{Cross-Reference:} For the full epistemological treatment, see Appendix UNMAPPED_MSC,
Section~\ref{sec:meta-framework-model}.


%%%%%%%%%%%%%%%%%%%%%%%%%%%%%%%%%%%%%%%%%%%%%%%%%%%%%%%%%%%%%%%%%%%%%%%%%%%%%%%
% SECTION 3: THE 10 ADVANTAGES OF EBF
%%%%%%%%%%%%%%%%%%%%%%%%%%%%%%%%%%%%%%%%%%%%%%%%%%%%%%%%%%%%%%%%%%%%%%%%%%%%%%%

\section{The 10 Advantages of EBF over Classical Frameworks}
\label{sec:consulting-10advantages}

\subsection{Overview}

\begin{table}[htbp]
\centering
\caption{10 Structural Advantages of EBF}
\label{tab:10-advantages}
\renewcommand{\arraystretch}{1.3}
\begin{tabular}{@{}clcc@{}}
\toprule
\textbf{\#} & \textbf{Advantage} & \textbf{Classical} & \textbf{EBF} \\
\midrule
1 & Output type & Descriptive & \textbf{Generative} \\
2 & Measurement & Qualitative & \textbf{Quantitative} \\
3 & Learning & Static & \textbf{Cumulative} \\
4 & Focus & Business structures & \textbf{Human behavior} \\
5 & Context & Ignored & \textbf{8$\Psi$ dimensions} \\
6 & Analysis level & Single & \textbf{Multi-level (9 levels)} \\
7 & Evidence base & Heuristic & \textbf{50+ years experimental} \\
8 & On failure & Terminal & \textbf{Informative} \\
9 & Technology & Manual only & \textbf{AI-enabled (Beatrix)} \\
10 & Competitive position & Commodity & \textbf{Cumulative advantage} \\
\bottomrule
\end{tabular}
\end{table}

\subsection{Detailed Analysis}

\paragraph{Advantage 1: Generative vs.\ Descriptive.}

\begin{center}
\begin{tabular}{@{}p{6.5cm}p{6.5cm}@{}}
\textbf{Classical (Five Forces):} & \textbf{EBF:} \\
\midrule
Output: ``Buyer power is high'' & Output: ``$\Psi_T$ = 82\% impact $\rightarrow$'' \\
$\rightarrow$ Now what? Consultant must figure out. & ``Weekly rewards increase participation by 12\% $\pm$3\%'' \\
\end{tabular}
\end{center}

\paragraph{Advantage 2: Quantitative vs.\ Qualitative.}

\begin{center}
\begin{tabular}{@{}p{6.5cm}p{6.5cm}@{}}
\textbf{Classical (BCG Matrix):} & \textbf{EBF:} \\
\midrule
``Product X is a Question Mark'' & ``Intervention X has 73\% success probability'' \\
What does that mean? 60\%? 30\%? 90\%? & ``(95\% CI: 64\%--82\%), ROI = 340\%'' \\
\end{tabular}
\end{center}

\paragraph{Advantage 3: Learning vs.\ Static.}

\begin{center}
\begin{tabular}{@{}p{6.5cm}p{6.5cm}@{}}
\textbf{Classical:} & \textbf{EBF:} \\
\midrule
Project 1: Recommendation $\rightarrow$ Success & Project 1: Test $\rightarrow$ Bayesian update \\
Project 2: Recommendation $\rightarrow$ Failure & Project 2: Prior already informed \\
Project 3: No improvement from P1, P2 & Project 100: Significantly better than P1 \\
\end{tabular}
\end{center}

\paragraph{Advantage 4: Behavioral vs.\ Structural Focus.}

Classical frameworks analyze \textit{markets and organizations}. EBF analyzes \textit{human behavior}:
\begin{itemize}[nosep]
    \item WHO: Who is the decision-maker?
    \item AWARE: What is psychologically salient?
    \item READY: Is the person willing to act?
    \item STAGE: Where in the behavioral change journey?
\end{itemize}

\paragraph{Advantage 5: Context-Sensitive vs.\ Context-Blind.}

\begin{center}
\begin{tabular}{@{}p{6.5cm}p{6.5cm}@{}}
\textbf{Classical:} & \textbf{EBF:} \\
\midrule
``Best practice: Opt-out defaults'' & ``Opt-out effect modulated by:'' \\
Works in Culture A, fails in Culture B. & $\Psi_{cultural}$, $\Psi_{trust}$ \\
Framework doesn't explain why. & ``For context X: +9\%, for context Y: +2\%'' \\
\end{tabular}
\end{center}

\paragraph{Advantage 6: Multi-Level vs.\ Single-Level.}

EBF's WHO dimension captures 9 aggregation levels:
\begin{center}
$L=0$ (Individual) $\rightarrow$ $L=1$ (Dyad) $\rightarrow$ $L=2$ (Family) $\rightarrow$ ... $\rightarrow$ $L=\infty$ (Meta)
\end{center}
The same decision has different answers depending on which level is salient.

\paragraph{Advantage 7: Experimentally Validated vs.\ Heuristic.}

\begin{center}
\begin{tabular}{@{}p{6.5cm}p{6.5cm}@{}}
\textbf{Classical:} & \textbf{EBF:} \\
\midrule
Based on: Porter's intuition (1979) & Based on: 50+ years experimental research \\
Industrial economics theory & Kahneman \& Tversky, Fehr \& Schmidt \\
No experimental validation & Thaler, 1000+ replicated experiments \\
\end{tabular}
\end{center}

\paragraph{Advantage 8: Failure-Informative vs.\ Failure-Terminal.}

\begin{center}
\begin{tabular}{@{}p{6.5cm}p{6.5cm}@{}}
\textbf{Classical:} & \textbf{EBF:} \\
\midrule
Recommendation fails: & Intervention fails: \\
``Framework doesn't fit here'' & ``$\Psi_T$ wasn't the bottleneck'' \\
Start over with different framework & Framework shows: Check $\Psi_I$ next \\
No learning & Systematic iteration \\
\end{tabular}
\end{center}

\paragraph{Advantage 9: Technology-Enabled vs.\ Manual-Only.}

\begin{center}
\begin{tabular}{@{}p{6.5cm}p{6.5cm}@{}}
\textbf{Classical:} & \textbf{EBF:} \\
\midrule
Five Forces: 5 fields & 10C $\times$ 8$\Psi$ $\times$ 6 FEPSDE = 384 dimensions \\
Fits in human working memory & Too complex for unaided cognition \\
Cannot be automated & \textbf{Beatrix}: AI-enabled analysis \\
Quality depends on consultant & Quality guaranteed by framework \\
\end{tabular}
\end{center}

\paragraph{Advantage 10: Cumulative Advantage vs.\ Commodity.}

\begin{center}
\begin{tabular}{@{}p{6.5cm}p{6.5cm}@{}}
\textbf{Classical:} & \textbf{EBF:} \\
\midrule
Five Forces is public knowledge & EBF + Beatrix + project data \\
Every MBA learns it & Proprietary, learning system \\
No differentiation & Each project improves predictions \\
$\rightarrow$ Commodity & $\rightarrow$ Cumulative moat \\
\end{tabular}
\end{center}


%%%%%%%%%%%%%%%%%%%%%%%%%%%%%%%%%%%%%%%%%%%%%%%%%%%%%%%%%%%%%%%%%%%%%%%%%%%%%%%
% SECTION 4: THE 20 KPIs
%%%%%%%%%%%%%%%%%%%%%%%%%%%%%%%%%%%%%%%%%%%%%%%%%%%%%%%%%%%%%%%%%%%%%%%%%%%%%%%

\section{The 20 KPIs: Quantified Comparison}
\label{sec:consulting-kpis}

\subsection{Part 1: 10 Consulting Firm KPIs}

\begin{table}[htbp]
\centering
\caption{Consulting Firm KPIs: Classical vs.\ EBF}
\label{tab:firm-kpis}
\renewcommand{\arraystretch}{1.3}
\small
\begin{tabular}{@{}clccc@{}}
\toprule
\textbf{\#} & \textbf{KPI} & \textbf{Classical} & \textbf{EBF} & \textbf{$\Delta$} \\
\midrule
1 & Revenue per Consultant & CHF 400k & CHF 650k & \textbf{+63\%} \\
2 & Project Margin & 35\% & 52\% & \textbf{+17pp} \\
3 & Utilization Rate & 65\% & 78\% & \textbf{+13pp} \\
4 & Client Retention & 40\% & 72\% & \textbf{+32pp} \\
5 & Time to Diagnosis & 4--6 weeks & 3--5 days & \textbf{-85\%} \\
6 & Knowledge Retention & 20\% & 85\% & \textbf{+65pp} \\
7 & Junior Effectiveness & 30\% of Senior & 75\% of Senior & \textbf{+45pp} \\
8 & Proposal Win Rate & 25\% & 45\% & \textbf{+20pp} \\
9 & Scalability (Projects/Partner) & 4 & 12 & \textbf{3$\times$} \\
10 & IP Value & Low (Commodity) & High (Proprietary) & \textbf{$\infty$} \\
\bottomrule
\end{tabular}
\end{table}

\paragraph{Key Drivers:}
\begin{itemize}[nosep]
    \item \textbf{Knowledge Retention (+65pp):} Knowledge in system, not heads
    \item \textbf{Junior Effectiveness (+45pp):} Beatrix enables junior output at senior level
    \item \textbf{Time to Diagnosis (-85\%):} Automated analysis vs.\ manual workshops
    \item \textbf{Scalability (3$\times$):} Technology leverage vs.\ human-only
\end{itemize}

\subsection{Part 2: 10 Client Value KPIs}

\begin{table}[htbp]
\centering
\caption{Client Value KPIs: Classical vs.\ EBF}
\label{tab:client-kpis}
\renewcommand{\arraystretch}{1.3}
\small
\begin{tabular}{@{}clccc@{}}
\toprule
\textbf{\#} & \textbf{KPI} & \textbf{Classical} & \textbf{EBF} & \textbf{$\Delta$} \\
\midrule
1 & ROI of Recommendation & 150\% & 450\% & \textbf{3$\times$} \\
2 & Implementation Success & 30\% & 75\% & \textbf{+45pp} \\
3 & Time to Results & 12--18 months & 3--6 months & \textbf{-70\%} \\
4 & Prediction Accuracy & N/A & 73\% & \textbf{$\infty$} \\
5 & Actionability & 40\% & 85\% & \textbf{+45pp} \\
6 & Knowledge Transfer & 15\% & 60\% & \textbf{+45pp} \\
7 & Sustainability (2 years) & 25\% & 65\% & \textbf{+40pp} \\
8 & Risk Reduction & Low & High & \textbf{Qual.} \\
9 & Cost per Outcome & CHF 50k/\% & CHF 8k/\% & \textbf{-84\%} \\
10 & Capability Building & 10\% & 55\% & \textbf{+45pp} \\
\bottomrule
\end{tabular}
\end{table}

\paragraph{Key Drivers:}
\begin{itemize}[nosep]
    \item \textbf{Prediction Accuracy ($\infty$):} Classical makes no predictions; EBF achieves 73\%
    \item \textbf{Cost per Outcome (-84\%):} Targeted interventions vs.\ trial-and-error
    \item \textbf{Implementation Success (+45pp):} Behavioral barriers identified upfront
    \item \textbf{ROI (3$\times$):} Right lever identified, not guessed
\end{itemize}

\subsection{The SwissHealth Validation}

\textit{Cross-Reference:} Full case in Appendix CTW, Section~\ref{sec:V-example-org}.

\begin{table}[htbp]
\centering
\caption{SwissHealth Case: KPI Comparison}
\label{tab:swisshealth-kpis}
\renewcommand{\arraystretch}{1.3}
\begin{tabular}{@{}lccc@{}}
\toprule
\textbf{KPI} & \textbf{Classical} & \textbf{EBF} & \textbf{$\Delta$} \\
\midrule
\multicolumn{4}{l}{\textit{Consulting Firm Perspective:}} \\
Time to Diagnosis & 6 weeks & 5 days & -88\% \\
Team Composition & 1 Partner + 2 Seniors & 1 Manager + 2 Juniors & -40\% cost \\
Project Margin & 28\% & 54\% & +26pp \\
\midrule
\multicolumn{4}{l}{\textit{Client Value Perspective:}} \\
Recommendation & ``Increase bonus'' & ``Redesign context'' & Qualitative \\
Prediction & ``Should help'' & ``+26\% $\pm$4\%'' & Quantified \\
Implementation Cost & CHF 1.6M/year & CHF 50k one-time & -97\% \\
Result & +1\% participation & +26\% participation & 26$\times$ \\
ROI & \textbf{Negative} & \textbf{+1,200\%} & $\infty$ \\
\bottomrule
\end{tabular}
\end{table}


%%%%%%%%%%%%%%%%%%%%%%%%%%%%%%%%%%%%%%%%%%%%%%%%%%%%%%%%%%%%%%%%%%%%%%%%%%%%%%%
% SECTION 5: BEATRIX OPERATIONALIZATION
%%%%%%%%%%%%%%%%%%%%%%%%%%%%%%%%%%%%%%%%%%%%%%%%%%%%%%%%%%%%%%%%%%%%%%%%%%%%%%%

\section{Beatrix: Technology-Enabled Behavioral Consulting}
\label{sec:consulting-beatrix}

\subsection{The Operationalization Challenge}

EBF's power comes with complexity:
\begin{equation}
\text{Analysis Space} = 8 \text{ (10C)} \times 8 \text{ (8$\Psi$)} \times 6 \text{ (FEPSDE)} = 384 \text{ dimensions}
\end{equation}

This exceeds human working memory capacity ($7 \pm 2$ items). Without technology,
EBF cannot be systematically applied.

\subsection{The Beatrix Architecture}

\begin{tcolorbox}[colback=purple!5!white, colframe=purple!75!black,
    title=\textbf{Beatrix: Framework-to-Model Generator}]

\begin{center}
\begin{tabular}{@{}c@{}}
\textbf{FRAMEWORK LAYER (EBF)} \\
10C CORE Questions $\times$ 8$\Psi$ Dimensions $\times$ FEPSDE \\
$\downarrow$ \\
\textbf{BEATRIX ANALYSIS ENGINE} \\
Diagnoses which dimensions are binding constraints \\
$\downarrow$ \\
\textbf{MODEL GENERATION} \\
Specific, testable predictions for this context \\
$\downarrow$ \\
\textbf{INTERVENTION DESIGN} \\
Actionable recommendations with predicted impact \\
\end{tabular}
\end{center}

\textbf{Key insight:} Beatrix uses the \textit{framework} to generate context-specific \textit{models}.
When a model fails, Beatrix doesn't start over---it queries the framework for the next hypothesis.

\end{tcolorbox}

\subsection{The Success Probability Mathematics}

\begin{tcolorbox}[colback=yellow!5!white, colframe=orange!75!black,
    title=\textbf{Model-Only vs.\ Framework + Beatrix}]

\textbf{Model-Only Consulting:}
\begin{verbatim}
IF model works -> success
IF model fails -> project fails (or restart)
Learning = 0 (each project independent)
\end{verbatim}

\textbf{Framework + Beatrix:}
\begin{verbatim}
IF model works -> success + update priors
IF model fails -> framework generates next model
Learning = cumulative (improves future projects)
\end{verbatim}

\textbf{Mathematical formulation:}
\begin{equation}
P(\text{success})_{\text{model-only}} = p_1 \approx 0.30
\end{equation}
\begin{equation}
P(\text{success})_{\text{framework}} = 1 - \prod_{i=1}^{n} (1 - p_i) \approx 1 - (1-p)^n
\end{equation}

With $n = 5$ iterations and $p = 0.30$ per model:
\begin{itemize}[nosep]
    \item Model-only: 30\% success rate
    \item Framework: $1 - (0.70)^5 = 83\%$ success rate
\end{itemize}

\end{tcolorbox}

\subsection{Without vs.\ With Beatrix}

\begin{table}[htbp]
\centering
\caption{Consulting Process: Without vs.\ With Beatrix}
\label{tab:beatrix-comparison}
\renewcommand{\arraystretch}{1.3}
\begin{tabular}{@{}p{3cm}p{5cm}p{5cm}@{}}
\toprule
\textbf{Step} & \textbf{Without Beatrix} & \textbf{With Beatrix} \\
\midrule
Diagnosis & Interviews, workshops, intuition (4--6 weeks) & Data $\rightarrow$ Beatrix $\rightarrow$ output (3--5 days) \\
Hypothesis & ``Bonus too low'' (intuition) & ``$\Psi_T$ = 82\% impact'' (data-driven) \\
Model & ``Increase bonus to CHF 400'' & ``Switch to weekly rewards'' \\
If fails & ``Incentives don't work here'' & Query framework: next $\Psi$-dimension \\
Learning & Ad hoc, lost when staff leaves & Systematic, accumulated in system \\
\bottomrule
\end{tabular}
\end{table}


%%%%%%%%%%%%%%%%%%%%%%%%%%%%%%%%%%%%%%%%%%%%%%%%%%%%%%%%%%%%%%%%%%%%%%%%%%%%%%%
% SECTION 6: IMPLICATIONS
%%%%%%%%%%%%%%%%%%%%%%%%%%%%%%%%%%%%%%%%%%%%%%%%%%%%%%%%%%%%%%%%%%%%%%%%%%%%%%%

\section{Implications for Consulting Practice}
\label{sec:consulting-implications}

\subsection{The Business Model Shift}

\begin{tcolorbox}[colback=gray!5!white, colframe=gray!75!black,
    title=\textbf{From Handicraft to Industrialization}]

\textbf{Classical Consulting Business Model:}
\begin{quote}
Sell TIME of smart people.\\
Problem: Time is finite, intelligence doesn't scale.\\
Result: High cost, variable quality.
\end{quote}

\textbf{EBF + Beatrix Business Model:}
\begin{quote}
Sell OUTCOMES from a system.\\
Advantage: System scales, learns, improves.\\
Result: Lower cost, consistent quality.
\end{quote}

\textbf{The difference is between handicraft and industrialization.}

Both have their place. But only one scales.

\end{tcolorbox}

\subsection{Implications by Stakeholder}

\paragraph{For Consulting Firms:}
\begin{itemize}[nosep]
    \item Invest in framework + technology, not just talent
    \item Knowledge retention becomes competitive advantage
    \item Junior leverage increases dramatically
    \item Differentiation through prediction, not description
\end{itemize}

\paragraph{For Clients:}
\begin{itemize}[nosep]
    \item Demand quantified predictions, not just recommendations
    \item Expect measurable ROI commitments
    \item Value knowledge transfer over dependency
    \item Select consultants by methodology, not just reputation
\end{itemize}

\paragraph{For the Industry:}
\begin{itemize}[nosep]
    \item Behavioral consulting becomes systematizable
    \item Technology creates new competitive dynamics
    \item ``Best practice'' gives way to ``best framework''
    \item Implementation success rates can dramatically improve
\end{itemize}


%%%%%%%%%%%%%%%%%%%%%%%%%%%%%%%%%%%%%%%%%%%%%%%%%%%%%%%%%%%%%%%%%%%%%%%%%%%%%%%
% SECTION 7: SUMMARY
%%%%%%%%%%%%%%%%%%%%%%%%%%%%%%%%%%%%%%%%%%%%%%%%%%%%%%%%%%%%%%%%%%%%%%%%%%%%%%%

\section{Summary}
\label{sec:consulting-summary}

\begin{tcolorbox}[colback=green!5!white, colframe=green!75!black,
    title=\textbf{DOMAIN-CONSULTING: Key Takeaways}]

\textbf{1. The Core Distinction:}
\begin{itemize}[nosep]
    \item Classical frameworks = descriptive templates (``what is'')
    \item EBF = generative framework (``what will be if...'')
\end{itemize}

\textbf{2. The 10 Advantages:}
\begin{enumerate}[nosep]
    \item Generative vs.\ descriptive
    \item Quantitative vs.\ qualitative
    \item Learning vs.\ static
    \item Behavioral vs.\ structural focus
    \item Context-sensitive vs.\ context-blind
    \item Multi-level vs.\ single-level
    \item Experimentally validated vs.\ heuristic
    \item Failure-informative vs.\ failure-terminal
    \item Technology-enabled vs.\ manual-only
    \item Cumulative advantage vs.\ commodity
\end{enumerate}

\textbf{3. The 20 KPIs:}
\begin{itemize}[nosep]
    \item 10 firm KPIs: Revenue +63\%, Knowledge retention +65pp, Scalability 3$\times$
    \item 10 client KPIs: ROI 3$\times$, Implementation success +45pp, Cost per outcome -84\%
\end{itemize}

\textbf{4. The Technology Enabler:}
\begin{itemize}[nosep]
    \item Beatrix operationalizes the 384-dimension analysis space
    \item Success probability: 30\% (model-only) $\rightarrow$ 83\% (framework)
\end{itemize}

\end{tcolorbox}

\begin{center}
\fbox{\parbox{0.9\textwidth}{
\centering
\textbf{The Bottom Line}\\[0.5em]
Classical consulting sells \textit{time and intuition}.\\
EBF consulting sells \textit{predictions and outcomes}.\\[0.5em]
The difference is measurable: 3$\times$ ROI, -84\% cost per outcome, +45pp success rate.
}}
\end{center}


%%%%%%%%%%%%%%%%%%%%%%%%%%%%%%%%%%%%%%%%%%%%%%%%%%%%%%%%%%%%%%%%%%%%%%%%%%%%%%%
% BACK MATTER
%%%%%%%%%%%%%%%%%%%%%%%%%%%%%%%%%%%%%%%%%%%%%%%%%%%%%%%%%%%%%%%%%%%%%%%%%%%%%%%

%------------------------------------------------------------------------------
% GLOSSARY SECTION
%------------------------------------------------------------------------------

\section{Glossary}
\label{sec:consulting-glossary}

Key terms used in this appendix follow the definitions in Appendix~G (Glossary):
\begin{itemize}[nosep]
    \item \textbf{Framework:} Organizing structure for analysis; generates models
    \item \textbf{Model:} Specific, testable prediction derived from framework
    \item \textbf{10C CORE:} Eight fundamental questions (WHO, WHAT, HOW, WHEN, WHERE, AWARE, READY, STAGE)
    \item \textbf{8$\Psi$:} Eight context dimensions that modulate behavior
    \item \textbf{Beatrix:} AI-assisted behavioral analysis system
    \item \textbf{KPI:} Key Performance Indicator
    \item \textbf{ROI:} Return on Investment
    \item \textbf{Schablone:} German for ``template''---a structured form with fields to fill
\end{itemize}

%------------------------------------------------------------------------------
% CRITICAL FOUNDATIONS
%------------------------------------------------------------------------------

\section{Critical Foundations}
\label{sec:consulting-foundations}

This appendix builds on:
\begin{enumerate}[nosep]
    \item \textbf{Appendix UNMAPPED_MSC:} Framework vs.\ Model epistemology
    \item \textbf{Appendix CTW:} Context ($\Psi$) dimensions and SwissHealth case
    \item \textbf{Appendix UNMAPPED_HOW:} Complementarity theory (10C CORE-HOW)
    \item \textbf{Appendix LLM1:} Beatrix technical methodology
\end{enumerate}

%------------------------------------------------------------------------------
% OPEN ISSUES
%------------------------------------------------------------------------------

\section{Open Issues and Future Research}
\label{sec:consulting-open}

\begin{enumerate}
\item \textbf{KPI Validation:} The 20 KPIs presented require systematic validation
across a larger sample of consulting engagements. Current estimates are based on
FehrAdvice project data (N$\approx$200 projects).

\item \textbf{Industry Variation:} Do the 10 advantages hold equally across industries
(healthcare, finance, retail, public sector)? Preliminary evidence suggests yes,
but systematic comparison is needed.

\item \textbf{Technology Dependency:} As Beatrix capabilities expand, the gap between
EBF and classical consulting may widen further. Longitudinal tracking is needed.

\item \textbf{Training Requirements:} What is the optimal training path for consultants
transitioning from classical to EBF-based methodology? Current estimate: 40--60 hours.

\item \textbf{Client Readiness:} Not all clients are ready for prediction-based consulting.
Segmentation of client readiness and appropriate positioning remains an open question.
\end{enumerate}

%------------------------------------------------------------------------------
% REFERENCES SECTION
%------------------------------------------------------------------------------

\section{References}
\label{sec:consulting-references}

For the complete bibliography, see \texttt{bibliography/bcm\_master.bib}.

Key references for this appendix:
\begin{itemize}[nosep]
    \item Porter, M.E. (1979). How competitive forces shape strategy. \textit{Harvard Business Review}.
    \item Waterman, R.H., Peters, T.J., \& Phillips, J.R. (1980). Structure is not organization. \textit{Business Horizons}.
    \item Kahneman, D. (2011). \textit{Thinking, Fast and Slow}. Farrar, Straus and Giroux.
    \item Thaler, R.H., \& Sunstein, C.R. (2008). \textit{Nudge}. Yale University Press.
    \item Fehr, E., \& Schmidt, K.M. (1999). A theory of fairness, competition, and cooperation. \textit{Quarterly Journal of Economics}.
\end{itemize}

% Master bibliography reference
\nocite{bcm_master}

%%%%%%%%%%%%%%%%%%%%%%%%%%%%%%%%%%%%%%%%%%%%%%%%%%%%%%%%%%%%%%%%%%%%%%%%%%%%%%%
% END OF APPENDIX DOMAIN-CONSULTING
%%%%%%%%%%%%%%%%%%%%%%%%%%%%%%%%%%%%%%%%%%%%%%%%%%%%%%%%%%%%%%%%%%%%%%%%%%%%%%%
